\documentclass[11pt]{article}
\usepackage[margin=1in]{geometry}
\usepackage[round]{natbib}
\usepackage{graphicx, booktabs, amsmath, hyperref}

% Find every file by name in the paper_output folder
\makeatletter\def\input@path{{paper_output/}}\makeatother
\graphicspath{{paper_output/}}
% \stat{name}: the number in name.txt
\newcommand{\stat}[1]{\input{#1.txt}\unskip}
% \fig{name}{caption}: name.pdf as a numbered figure
\newcommand{\fig}[2]{\begin{figure}[htbp]\centering
  \includegraphics[width=0.8\textwidth]{#1.pdf}\caption{#2}\label{fig:#1}\end{figure}}
% \tbl{name}{caption}: name.tex as a numbered table
\newcommand{\tbl}[2]{\begin{table}[htbp]\centering
  \caption{#2}\label{tbl:#1}\input{#1.tex}\end{table}}

\title{Prices, Promotions, and Quantities in a Synthetic Retail Panel}
\author{A. Student}

\begin{document}
\maketitle

\section{Introduction}

Building on classic models of consumer search \citep{stigler1961economics}
and sales \citep{varian1980model}, this paper studies how posted prices
relate to quantities sold in a synthetic retail panel.

\section{Data}

The panel has \stat{n-observations} observations, with an average posted
price of \$\stat{mean-price}.

Table \ref{tbl:results-by-product} summarizes the panel by product.

\tbl{results-by-product}{Mean price (\$), mean quantity, and promotion rate by product.}

\section{Results}

Figure \ref{fig:price-quantity} plots quantity against posted price.

\fig{price-quantity}{Quantity sold against posted price, with an ordinary least squares fit.}

A one-unit increase in price is associated with a change of
\stat{price-quantity-slope} units sold, in the spirit of \citet{stigler1961economics}.

\bibliographystyle{plainnat}
\bibliography{references}

\end{document}
